% common/tex/slides.tex
% Copyright (C) 2026 Hirokazu Maesaka (RIKEN SPring-8 Center)
% SPDX-License-Identifier: CC-BY-4.0
%
% Slide markers. Each marker writes one line "kind:payload" to build/<topic>/<jobname>.slides during compilation.
% common/slides/build_slides.py reads that file and builds the PowerPoint deck. The markers print nothing in the PDF.
% Arguments are expanded at the point of use, so \en{...} and \ja{...} select the language of the current build.
% Do not use fragile commands (for example \cite or \label) inside slide arguments.

\newwrite\slidesfile
\AtBeginDocument{\immediate\openout\slidesfile=\jobname.slides\relax}
\AtEndDocument{\immediate\closeout\slidesfile}

% #1: entry kind, #2: payload. The payload is expanded, then written as plain characters.
\newcommand{\slideentry}[2]{%
  \begingroup
  \edef\slidetmp{#2}%
  \edef\slidedet{\detokenize\expandafter{\slidetmp}}%
  \immediate\write\slidesfile{#1:\slidedet}%
  \endgroup}

% Raw entries are written without expansion. Use them for math and paths, which contain fragile commands.
\newcommand{\slideentryraw}[2]{%
  \immediate\write\slidesfile{#1:\detokenize{#2}}}

\newcommand{\slidetitle}[1]{\slideentry{title}{#1}}
\newcommand{\slidepoint}[1]{\slideentry{point}{#1}}
\newcommand{\slidemath}[1]{\slideentryraw{math}{#1}}
\newcommand{\slidefigure}[1]{\slideentryraw{figure}{#1}}
\newcommand{\slidenote}[1]{\slideentry{note}{#1}}
